% Part II -- Python Foundations
% Chapter 3: Organizing Data

\h{Organizing Data}
\label{ch:03}

A single variable holds one value. Many tasks require several related values:
room names, floor areas, material codes, or measurements. Python provides
collections for storing these values together. This chapter introduces four
collections and explains when each one is useful.

\begin{NoteBox}
\textbf{Prerequisites.} You should understand values, variables, assignment,
basic types, and \c{print()}.
\end{NoteBox}

\begin{tcolorbox}[title=Learning outcomes,colframe=orange,colback=white,arc=2mm]
After completing this chapter, you should be able to:
\begin{itemize}
    \item create, read, and modify a list;
    \item store labeled values in a dictionary;
    \item use a tuple for a fixed record and a set for unique values;
    \item test whether a collection contains a value; and
    \item choose a suitable collection for a simple situation.
\end{itemize}
\end{tcolorbox}

\hh{Why collections are useful}

Without a collection, four room names require four variables.

\begin{lstlisting}[
    language=python,
    style=block,
    caption={Four separate variables},
    label={c03_separate_names},
    captionpos=b
]
room_1 = "Lobby"
room_2 = "Studio"
room_3 = "Office"
room_4 = "Cafe"
\end{lstlisting}

A list stores the same information under one name.

\begin{lstlisting}[
    language=python,
    style=block,
    caption={Store several room names in one list},
    label={c03_first_list},
    captionpos=b
]
room_names = ["Lobby", "Studio", "Office", "Cafe"]
print(room_names)
\end{lstlisting}

The square brackets identify a list. Commas separate its items.

\begin{ExplainBox}{A collection encodes a relationship among values}
Putting values together is not only a storage convenience. The collection
type states what relationship the program relies on. A list says that order
and position matter. A dictionary says that each value is retrieved through a
meaningful key. A tuple says that a small ordered record should not be changed
in place. A set says that membership and uniqueness matter more than position.
Choosing the wrong relationship produces awkward code even when the values
themselves are correct.
\end{ExplainBox}

Collections also have behavioral properties:
\begin{itemize}
    \item \textbf{ordered} means iteration follows a defined sequence;
    \item \textbf{mutable} means the existing object can be changed in place;
    \item \textbf{indexed} means a position such as \c{0} selects an item;
    \item \textbf{keyed} means a meaningful key selects a value;
    \item \textbf{unique} means duplicate members collapse into one member.
\end{itemize}

These properties are independent questions. A list is ordered, mutable, and
indexed, while a set is mutable and unique but not position-indexed. Chapter 9
will show why mutability also affects assignment and copying.

\hh{Create lists}

A list is an ordered, changeable collection. Its items may be numbers, strings,
booleans, or other values.

\begin{lstlisting}[
    language=python,
    style=block,
    caption={Create four lists},
    label={c03_list_examples},
    captionpos=b
]
floor_numbers = [1, 2, 3, 4]
room_areas_m2 = [24.0, 35.5, 18.0]
status_labels = ["Open", "Closed", "Review"]
empty_list = []

print(floor_numbers)
print(room_areas_m2)
print(status_labels)
print(empty_list)
\end{lstlisting}

Use \c{len()} to find the number of items.

\begin{lstlisting}[
    language=python,
    style=block,
    caption={Count list items},
    label={c03_list_length},
    captionpos=b
]
room_names = ["Lobby", "Studio", "Office", "Cafe"]
print(len(room_names))
\end{lstlisting}

\hh{Read list items by index}

Each item has a position called an \emph{index}. Python begins counting at
zero.

\begin{lstlisting}[
    language=python,
    style=block,
    caption={Read list items by position},
    label={c03_list_index},
    captionpos=b
]
room_names = ["Lobby", "Studio", "Office", "Cafe"]

print(room_names[0])
print(room_names[1])
print(room_names[-1])
\end{lstlisting}

Expected output:

\begin{lstlisting}[
    language=text,
    style=block,
    caption={List indexing output},
    label={c03_list_index_output},
    captionpos=b
]
Lobby
Studio
Cafe
\end{lstlisting}

An index of \c{-1} means the last item.

\begin{NoteBox}
\textbf{Stop and predict.} What will \c{room_names[2]} display? What is the
largest valid positive index for this four-item list?
\end{NoteBox}

An index outside the list produces an \c{IndexError}.

\hh{Select part of a list}

A slice selects a range of items. The stop position is not included.

\begin{lstlisting}[
    language=python,
    style=block,
    caption={Select list slices},
    label={c03_list_slice},
    captionpos=b
]
levels = ["B1", "L1", "L2", "L3", "Roof"]

print(levels[1:4])
print(levels[:2])
print(levels[3:])
\end{lstlisting}

Expected output:

\begin{lstlisting}[
    language=text,
    style=block,
    caption={Three list slices},
    label={c03_list_slice_output},
    captionpos=b
]
['L1', 'L2', 'L3']
['B1', 'L1']
['L3', 'Roof']
\end{lstlisting}

\begin{SyntaxBox}{A slice describes a half-open range}
Lists, strings, and tuples are all sequences -- ordered collections that
support this same slicing pattern. The pattern is \c{sequence[start:stop]}.
Python includes the item at \c{start} and stops before \c{stop}. Omitting
\c{start} means ``from the beginning''; omitting \c{stop} means ``through the
end.''
\end{SyntaxBox}

\bookfigure[0.88\textwidth]{Figures/List/List_operator.pdf}
{Visual guide to list indexing, slicing, and common list operations.}
{fig:ch03-list-operations}

\hh{Change a list}

Lists are mutable, which means their contents can change.

\begin{lstlisting}[
    language=python,
    style=block,
    caption={Replace one item and append another},
    label={c03_modify_list},
    captionpos=b
]
room_names = ["Lobby", "Studio", "Office"]

room_names[2] = "Meeting Room"
room_names.append("Cafe")

print(room_names)
\end{lstlisting}

Common list methods include:
\begin{itemize}
    \item \c{append(value)}: add one item at the end;
    \item \c{insert(index, value)}: add an item at a chosen position;
    \item \c{remove(value)}: remove the first matching value;
    \item \c{pop()}: remove and return the last item; and
    \item \c{clear()}: remove all items.
\end{itemize}

\begin{lstlisting}[
    language=python,
    style=block,
    caption={Insert, remove, and pop list items},
    label={c03_list_methods},
    captionpos=b
]
materials = ["Concrete", "Steel"]

materials.insert(1, "Timber")
materials.remove("Steel")
removed_item = materials.pop()

print(materials)
print(removed_item)
\end{lstlisting}

After \c{remove("Steel")}, the list contains \c{"Concrete"} and
\c{"Timber"}. Then \c{pop()} removes and returns \c{"Timber"}.

\hh{Check membership}

The \c{in} operator asks whether a collection contains a value. The result is
\c{True} or \c{False}.

\begin{lstlisting}[
    language=python,
    style=block,
    caption={Check list membership},
    label={c03_list_membership},
    captionpos=b
]
materials = ["Concrete", "Timber", "Glass"]

print("Timber" in materials)
print("Brick" in materials)
\end{lstlisting}

Membership checks are case-sensitive. \c{"timber"} and \c{"Timber"} are
different strings.

\hh{Create dictionaries}

A dictionary stores key--value pairs. Each key is a label used to retrieve its
value. Curly braces identify a dictionary, and a colon separates each key from
its value.

\begin{lstlisting}[
    language=python,
    style=block,
    caption={A dictionary describing one room},
    label={c03_first_dictionary},
    captionpos=b
]
room = {
    "name": "Studio",
    "area_m2": 48.0,
    "level": 2
}

print(room)
\end{lstlisting}

Use a key inside square brackets to read a value.

\begin{lstlisting}[
    language=python,
    style=block,
    caption={Read dictionary values by key},
    label={c03_dictionary_access},
    captionpos=b
]
room = {
    "name": "Studio",
    "area_m2": 48.0,
    "level": 2
}

print(room["name"])
print(room["area_m2"])
\end{lstlisting}

Unlike list indexes, dictionary keys communicate what the values mean.

\bookfigure[0.86\textwidth]{Figures/Dictionary/Dictionary.pdf}
{A dictionary connects descriptive keys to values rather than relying on numerical positions.}
{fig:ch03-dictionary-model}

\begin{TryBox}{Choose a collection before writing code}
For each case, choose a list, dictionary, tuple, or set and justify the choice:
an ordered room schedule; one room's named properties; an immutable RGB color;
and the unique material names found in a model. Then create one small example
of each collection in Python.
\end{TryBox}

\hh{Update a dictionary}

Assign to an existing key to replace its value. Assign to a new key to add a
new pair.

\begin{lstlisting}[
    language=python,
    style=block,
    caption={Replace one value and add another},
    label={c03_dictionary_update},
    captionpos=b
]
room = {
    "name": "Studio",
    "area_m2": 48.0
}

room["area_m2"] = 50.5
room["level"] = 2

print(room)
\end{lstlisting}

Use \c{pop(key)} to remove a pair and return its value.

\begin{lstlisting}[
    language=python,
    style=block,
    caption={Remove a dictionary pair},
    label={c03_dictionary_pop},
    captionpos=b
]
room = {
    "name": "Studio",
    "area_m2": 50.5,
    "level": 2
}

removed_level = room.pop("level")
print(removed_level)
print(room)
\end{lstlisting}

\hh{Use safe dictionary lookup}

Square-bracket lookup produces a \c{KeyError} when the requested key is
missing. The \c{get()} method can return a default value instead.

\begin{lstlisting}[
    language=python,
    style=block,
    caption={Use get for a key that may be missing},
    label={c03_dictionary_get},
    captionpos=b
]
room = {
    "name": "Studio",
    "area_m2": 50.5
}

print(room.get("name"))
print(room.get("occupancy", "Not recorded"))
\end{lstlisting}

Use square brackets when a key must exist. Use \c{get()} when a missing value
has a reasonable default.

Dictionary membership checks keys.

\begin{lstlisting}[
    language=python,
    style=block,
    caption={Check whether dictionary keys exist},
    label={c03_dictionary_membership},
    captionpos=b
]
room = {"name": "Studio", "area_m2": 50.5}

print("name" in room)
print("level" in room)
\end{lstlisting}

\hh{Use tuples for fixed records}

A tuple is an ordered collection that cannot be changed after creation.
Parentheses identify a tuple.

\begin{lstlisting}[
    language=python,
    style=block,
    caption={A fixed two-coordinate record},
    label={c03_first_tuple},
    captionpos=b
]
point_xy = (12.5, 8.0)

print(point_xy[0])
print(point_xy[1])
\end{lstlisting}

A tuple is useful when the number and meaning of its positions should remain
fixed. Use a list when items need to be added, removed, or replaced.

\hh{Use sets for unique values}

A set stores unique values. Repeated values are removed automatically. Sets do
not provide list-style positional indexing.

\begin{lstlisting}[
    language=python,
    style=block,
    caption={A set removes a repeated value},
    label={c03_first_set},
    captionpos=b
]
materials = {"Steel", "Glass", "Steel", "Concrete"}
print(materials)
\end{lstlisting}

The displayed order may vary. Use \c{set()} to create an empty set because
empty curly braces create an empty dictionary.

\begin{lstlisting}[
    language=python,
    style=block,
    caption={Add an item and check set membership},
    label={c03_set_membership},
    captionpos=b
]
materials = {"Steel", "Glass", "Concrete"}

materials.add("Timber")
print("Timber" in materials)
print(len(materials))
\end{lstlisting}

\hh{Choose a collection}

\begin{center}
\begin{tabular}{lll}
\textbf{Collection} & \textbf{Main feature} & \textbf{Use when} \\
List & Ordered and changeable & Position and sequence matter \\
Dictionary & Labeled key--value pairs & Values need meaningful labels \\
Tuple & Ordered and fixed & A small record should not change \\
Set & Unique values & Duplicates should be removed \\
\end{tabular}
\end{center}

Ask two questions before choosing: Does order matter? Does each value need a
label? These questions usually narrow the choice quickly.

Then test the choice against the operations the program must perform:
\begin{enumerate}
    \item Must duplicate values be preserved? Use a list or tuple, not a set.
    \item Must an item be retrieved by a descriptive name? Use a dictionary.
    \item Must insertion, deletion, or replacement occur in place? Use a
          mutable list, dictionary, or set.
    \item Is the object a fixed coordinate or record whose positions have
          stable meanings? A tuple may communicate that constraint.
    \item Is the central question ``has this value appeared?'' A set expresses
          that membership question directly.
\end{enumerate}

For example, room names in walking order belong in a list; room areas looked
up by room name belong in a dictionary; a fixed \c{(x, y)} coordinate can be a
tuple; and the unique material categories present in a dataset belong in a
set. The values may look similar, but the required relationships differ.

\hh{Common mistakes}

\begin{itemize}
    \item \textbf{Starting at index 1}: the first list index is 0.
    \item \textbf{Using a missing key}: check the key or use \c{get()}.
    \item \textbf{Expecting a set to preserve position}: sets are not indexed sequences.
    \item \textbf{Creating an empty set with curly braces}: use \c{set()}.
    \item \textbf{Changing a tuple}: create a new tuple instead.
\end{itemize}

\hh{Practice ladder}

\hhh{Level 0 -- Read}

\begin{lstlisting}[
    language=python,
    style=block,
    caption={Predict two list lookups},
    label={c03_practice_level0},
    captionpos=b
]
levels = ["B1", "L1", "L2"]
print(levels[0])
print(levels[-1])
\end{lstlisting}

\hhh{Level 1 -- Modify}

Add two room names to the list, replace the first item, and print the result.

\begin{lstlisting}[
    language=python,
    style=block,
    caption={Starting list for modification},
    label={c03_practice_level1},
    captionpos=b
]
rooms = ["Lobby", "Studio"]
print(rooms)
\end{lstlisting}

\hhh{Level 2 -- Complete}

Replace each \c{...} with a value to complete a dictionary that stores a room
name, area, and level. Then display only the name and area.

\begin{lstlisting}[
    language=python,
    style=block,
    caption={Complete three dictionary values},
    label={c03_practice_level2},
    captionpos=b
]
room = {
    "name": ...,
    "area_m2": ...,
    "level": ...
}

print(room["name"])
print(room["area_m2"])
\end{lstlisting}

\hhh{Level 3 -- Apply}

Create:
\begin{enumerate}
    \item a list containing four material names;
    \item a dictionary containing one material's name and unit cost; and
    \item a set created from a list that contains at least one duplicate.
\end{enumerate}
Display each collection and its length.

\textbf{Hints}
\begin{itemize}
    \item \textbf{Hint 1}: Use square brackets for the list.
    \item \textbf{Hint 2}: Use descriptive string keys in the dictionary.
    \item \textbf{Hint 3}: Pass the duplicate-containing list to \c{set()}.
\end{itemize}

\hhh{Level 4 -- Challenge}

Create a dictionary containing a room's name and a list of three finish
materials. Add a fourth material to the nested list and display the result.

\hh{Practice solutions}

\textbf{Level 0:} The output is \c{B1}, then \c{L2}.

\textbf{Level 2}

\begin{lstlisting}[
    language=python,
    style=block,
    caption={One solution to Level 2},
    label={c03_solution_level2},
    captionpos=b
]
room = {
    "name": "Seminar Room",
    "area_m2": 32.0,
    "level": 3
}

print(room["name"])
print(room["area_m2"])
\end{lstlisting}

\textbf{Level 3 -- one possible solution}

\begin{lstlisting}[
    language=python,
    style=block,
    caption={One solution to Level 3},
    label={c03_solution_level3},
    captionpos=b
]
materials = ["Concrete", "Steel", "Glass", "Timber"]
material_cost = {"name": "Concrete", "cost_per_m2": 42.0}

repeated_names = ["Steel", "Glass", "Steel", "Timber"]
unique_materials = set(repeated_names)

print(materials, len(materials))
print(material_cost, len(material_cost))
print(unique_materials, len(unique_materials))
\end{lstlisting}

\textbf{Level 4}

\begin{lstlisting}[
    language=python,
    style=block,
    caption={One solution to Level 4},
    label={c03_solution_level4},
    captionpos=b
]
room = {
    "name": "Studio",
    "finishes": ["Concrete", "Paint", "Timber"]
}

room["finishes"].append("Glass")
print(room)
\end{lstlisting}

\hh{Chapter checkpoint}

\begin{enumerate}
    \item Which brackets identify a list?
    \item What is the index of the first list item?
    \item What does a negative index of \c{-1} select?
    \item Does a slice include its stop position?
    \item Which method adds one item to the end of a list?
    \item What two parts make up a dictionary pair?
    \item When is \c{get()} useful?
    \item Why might you choose a tuple instead of a list?
    \item What happens to duplicate values in a set?
    \item How do you create an empty set?
\end{enumerate}

\hh{Checkpoint answers}

\begin{enumerate}
    \item Square brackets.
    \item Zero.
    \item The last item.
    \item No.
    \item \c{append()}.
    \item A key and a value.
    \item When a key may be missing and a default is appropriate.
    \item A tuple communicates that the record should remain fixed.
    \item They are removed.
    \item Use \c{set()}.
\end{enumerate}

\begin{tcolorbox}[title=Chapter summary,colframe=orange,colback=white,arc=2mm]
Lists store ordered, changeable items. Dictionaries connect labels to values.
Tuples are useful for fixed records, and sets store unique values. Chapter 4
uses boolean expressions and conditions to make decisions with these values.
\end{tcolorbox}

